\documentclass[english, parskip=full]{scrartcl}
\usepackage[english]{babel}  % Sprache auf Deutsch 
\usepackage[utf8]{inputenc}  % Kodierung der Datei
\usepackage[left=2cm, right=2cm, top=2cm, bottom=3cm, footskip=1.5cm]{geometry}
\usepackage{color}
\usepackage{graphicx, gensymb}
\usepackage{amsfonts, cancel, amssymb}
\usepackage{amsmath, amsthm, array, esint}
\usepackage{bm} % bold math symbols, e.g. \bm{\sigma} etc.
\usepackage{booktabs, multirow}  % schönere Tabellen
%\usepackage{scrlayer-scrpage}    % Manipulation des Seitenstils
%\pagestyle{scrheadings}  % Stil für die Seite setzen
\usepackage{tabularx}
\usepackage{setspace}
%\automark{section}  % Abschnittsnamen als Seitenbeschriftung 
%\setheadsepline{.5pt}    % Kopzeile durch Linie abtrennen
\usepackage[hidelinks]{hyperref}  % Links und weitere PDF-Features
\usepackage{typearea, tikz, pgffor, verbatim}
\usetikzlibrary{calc, arrows.meta,arrows, graphs, quotes, positioning, angles, babel}
\usepackage[cache=false, outputdir=build]{minted}
\usemintedstyle{vs}
\usepackage{placeins} % provides \FloatBarrier

\definecolor{bg}{rgb}{0.9,0.9,0.9}
\newcommand\numberthis{\addtocounter{equation}{1}\tag{\theequation}}
\usepackage{svg}
\usepackage{siunitx}
\usepackage{physics}
\usepackage{tensor}
\usepackage{slashed}
\usepackage{bbm}
\usepackage[compat=1.1.0]{tikz-feynman}

\usepackage{caption}
\captionsetup{
    % width=.8\textwidth,
    % indention=1cm,
    margin=0.5cm,
    format=plain,
    labelfont=bf
}
%\usepackage{subcaption}
%\subcaptionsetup{
%    indention=0cm,
%    format=hang,
%    margin=0.5cm,
%    labelfont=normalfont
%}
\usepackage[english,capitalise]{cleveref}
\usepackage[
    backend=biber,
    sorting=none,
    style=phys,
    giveninits=true,
    sortcites=true,
    citestyle=numeric-comp,
    % url=false,
    % doi=false,
    biblabel=brackets,
    % eprint=false,
    maxcitenames=3]{biblatex}
\usepackage{csquotes}
%\addbibresource{bibliography.bib}

\usepackage{mathtools}
% use \abs for absolute values
% use \abs for \left ... \right version and \abs[\bigg for \bigg version etc.
\let\abs\undefined
\let\bra\undefined
\let\braket\undefined
\DeclarePairedDelimiter\abs{\vert}{\vert}
\DeclarePairedDelimiter\kla{(}{)}
\DeclarePairedDelimiter\klb{[}{]}
\DeclarePairedDelimiter\klc{\{}{\}}
\DeclarePairedDelimiter\bra{\langle}{\rangle}
\DeclarePairedDelimiterX\brb[2]{\langle #1}{\rangle}{\delimsize\vert #2 }
\DeclarePairedDelimiterX\brc[3]{\langle #1}{\rangle}{\delimsize\vert #2 \delimsize\vert #3 }

\renewcommand{\i}{\text{i}}
\newcommand{\e}{\text{e}}
\DeclareMathOperator{\sgn}{sgn}
\newcommand{\kom}[1]{\overset{\substack{#1 \\ |}}{=}}
\newcommand{\koml}[2]{\overset{\substack{#1 \\ |}}{#2}}
\newcommand{\intx}[4]{\int_{#1}^{#2} #3 \,\mathrm{d}#4}
\newcommand{\intr}[2]{\int_{\mathbb{R}} #1 \, \mathrm{d}#2}
\newcommand{\intrnd}[3]{\int_{\mathbb{R}^{#1}} #2 \, \mathrm{d}^{#1}#3}
\newcommand{\intrpi}[2]{\int_{\mathbb{R}} #1 \, \frac{\mathrm{d}#2}{2\pi}}
\newcommand{\intrndpi}[3]{\int_{\mathbb{R}^{#1}} #2 \, \frac{\mathrm{d}^{#1}#3}{(2\pi)^{#1}}}
\newcommand{\oper}[2]{\vert #1 \rangle \langle #2 \vert}
\newcommand{\Text}[1]{\text{#1}}    % this is relevant because \text does not autocomplete to the default '\text{@}' but rather '\textbf' etc.
\newcommand{\powe}[1]{\e^{#1}}
\newcommand{\pow}[2]{{#1}^{#2}}
\newcommand{\Ket}[1]{\vert #1 \rangle}
\newcommand{\Bra}[1]{\langle #1 \vert}
\newcommand*{\Fourier}{\textsc{Fourier}\ }
\newcommand*{\F}{\mathcal{F}}
\DeclareMathOperator{\arsinh}{arsinh}
\DeclareMathOperator{\arcosh}{arcosh}
\DeclareMathOperator{\artanh}{artanh}
\DeclareMathOperator{\sinc}{sinc}
\DeclareMathOperator{\cscc}{cscc}
\DeclareMathOperator{\sinhc}{sinhc}
\DeclareMathOperator{\cschc}{cschc}
\newcommand{\conv}{\mathop{\scalebox{3.5}{\raisebox{-0.38ex}{\makebox[0.5\width][c]{$\ast$}}}}\limits}
\newcommand*{\unity}{\text{\usefont{U}{bbold}{m}{n}1}}   % operator one from :: https://tex.stackexchange.com/questions/566780/import-mathbb1-character-from-the-unicode-math-package

\renewcommand{\d}{\text{d}}
\renewcommand{\r}{\text{r}}
\renewcommand{\t}{\text{t}}
\renewcommand{\a}{\text{a}}
\renewcommand{\o}{\text{o}}
\renewcommand{\F}{\text{F}}
\renewcommand{\c}{\text{c}}
\newcommand{\A}{\text{A}}
\newcommand{\B}{\text{B}}
\newcommand{\R}{\text{R}}
\newcommand{\M}{\text{M}}
\newcommand{\m}{\text{m}}
\newcommand{\s}{\text{s}}
\newcommand{\f}{\text{f}}
\newcommand{\K}{\text{K}}
\newcommand{\hc}{\text{h.c.}}
\newcommand{\kf}{k_\F}
\newcommand{\vf}{v_\F}
\newcommand{\const}{\text{const.}}
\renewcommand{\L}{\text{L}}
\newcommand{\gray}[1]{\bgroup\color{white!40!black}#1\egroup}
\newcommand{\TODO}[1]{\bgroup\color{red!60!black}#1\egroup}
\newcommand\QnA[1]{\bgroup\color{blue}#1\egroup}
\newcommand\highlight[1]{\bgroup\color{green!60!black}#1\egroup}

\newcommand{\cabx}[3]{\begin{smallmatrix}\gray{A}&\gray{:}&#1 \\ \gray{B}&\gray{:}&#2 \\ \gray{\Xi}&\gray{:}&#3\end{smallmatrix}}
\newcommand{\zabx}[3]{\begin{smallmatrix}\gray{A}&\gray{:}&#1 \\ \gray{B}&\gray{:}&#2 \\ \gray{Z}&\gray{:}&#3\end{smallmatrix}}
\newcommand{\abx}[3]{\begin{smallmatrix}#1 \\ #2 \\ #3\end{smallmatrix}}

\newcommand{\normord}[1]{
  {:\mathrel{\mspace{1mu}#1\mspace{1mu}}:}
}
\newcommand{\varnormord}[1]{
  {\mathrel{\mspace{1mu}\substack{\ast\\[-1.5pt] \ast}#1\substack{\ast\\[-1.5pt] \ast}\mspace{1mu}}}
}

% punctuation styles (see https://tex.stackexchange.com/questions/56063/how-to-properly-typeset-footnotes-superscripts-after-punctuation-marks)
\let\origfootnote\footnote
\renewcommand{\footnote}[1]{\kern.06em\origfootnote{#1}}
\newcommand{\punctfootnote}[1]{\kern-.06em\origfootnote{#1}}

% Multiline equations
\newcommand{\bsplit}[1]{\begin{aligned}[t]#1\end{aligned}}

\newcommand*{\titel}{Operator Product Expansion}
\newcommand*{\autor}{Joachim Schwardt}

\hypersetup{pdfauthor={\autor}, pdftitle={\titel}}

%\titlehead{titlehead}
\title{\titel}
\author{\autor}
\date{\today}

\begin{document}

%\begin{titlepage}
\maketitle  % Titel setzen
%\tableofcontents  % Inhaltsverzeichnis setzen
%\end{titlepage}

\section{General Theory}
In this section we want to present the RG approach based on the Operator Product Expansion (OPE) following Fradkin \TODO{[ref book and lecture notes; slight differences between each]}.
There are a several ways one can interpret and use this RG approach.
The fundamental starting point here will always be to introduce some notion of a short-distance cut-off $\alpha$ that regulates the UV divergences and can be thought of as generating a finite band-width $\Lambda\sim \frac{1}{\alpha}$.
Starting from an initial $\alpha_0$ that might for example correspond to the validity regime of the bare theory one then seeks to increase $\alpha\rightarrow \alpha + \d \alpha$ while keeping the partition function constant.
The resulting RG flow may then be terminated by some minimal bandwidth, e.g. thermal fluctuations $\Lambda_\text{thermal} = \frac{1}{\beta v}$ with some typical velocity $v$ and the inverse temperature $\beta=\frac{1}{T}$.
\subsection{Perturbative Renormalization Group and OPE}
With the previous discussion in mind let us now consider a set of operators $\mathcal{O}$ that are functions of some real scalar field $\phi$ in $D$-dimensional space-time ; later we will always use $D=2$.
Then the scaling dimension $\Delta_{\mathcal{O}}$ of an operator is defined by the asymptotic relation
\begin{align}
\bra*{\mathcal{O}(x_1) \mathcal{O}(x_2)} \overset{x_1\rightarrow x_2}{\sim} \frac{C_{\mathcal{O}}^2}{|x_1-x_2|^{2\Delta_{\mathcal{O}}}}
\end{align}
where $C_\mathcal{O}$ is some normalization of the operator.
We can then define a general action via
\begin{align}
S[\phi] &= S_0[\phi] + \sum_j g_j \alpha^{\Delta_{j} - D} \int \mathcal{O}_j(x)\,\d^Dx
\end{align}
where $g_j$ are by construction dimensionless couplings and $S_0$ is an action corresponding to a fixed point (e.g. a kinetic term $(\nabla\phi)^2$).
The partition function can then be written as
\begin{align}
Z &= \tr \klb*{\powe{S[\phi]}} = \tr \klb*{\powe{-S_0[\phi]} \powe{\sum_j g_j \alpha^{\Delta_{j} - D} \int \mathcal{O}_j(x)\,\d^Dx}}
\end{align}
where the trace runs over all configurations and may for example be given by a path integral.
Defining the expectation value $\bra*{\mathcal{O}} = \frac{1}{Z_0} \tr \klb*{ \powe{S_0} \mathcal{O}}$ with respect to the fixed point action $S_0$ allows us to formally rewrite this as
\begin{align}
\frac{Z}{Z_0} &=1 + \sum_j \frac{g_j}{\alpha^{D-\Delta_{j}}} \int \bra*{\mathcal{O}_j(x)}\,\d^Dx + \frac{1}{2} \sum_{j,k} \frac{g_j g_k}{\alpha^{2D-\Delta_{j}-\Delta_{k}}} \int \bra*{\mathcal{O}_j(x_1) \mathcal{O}_k(x_2)}\,\d^Dx_1\,\d^Dx_2 + \dots 
\end{align}
where $Z_0=\tr \powe{S_0}$ and we expanded the exponential containing the perturbations to the action $S_0$.\\
All of the equations here contain the short-distance cut-off $\alpha$ we alluded to in the introduction -- but how does it enter into the calculation?
The answer to this is that we restrict all of the integrations in the series expansion such that $|x_j - x_k| > \alpha$ for any combination of coordinates.
To implement this idea we now write $\alpha'=s\alpha$ where $s=\powe{\d l} = 1 + \d l$ describes an infinitesimal change and try to adjust the couplings $g_j$ such that the new theory looks structurally identical.
Equivalently we have $\alpha(l)=\alpha_0\powe{l}$ and the condition $\frac{\d}{\d l} \frac{Z}{Z_0} = 0$.
Let us see what happens when we differentiate the $N$-th order term in the expansion.
For that purpose we will introduce the multi-index notations $j = (j_1,\dots,j_N)$, $g!=g_{j_1} \dots g_{j_N}$ and $|g|=g_{j_1} + \dots + g_{j_N}$ which allows us to write $\frac{Z}{Z_0} = \sum_{N\ge 0} \frac{1}{N!} Z_N$ where
\begin{align}
Z_N &= \sum_j \frac{g!}{\alpha^{ND - |\Delta_j|}} I_N(j) = \sum_j \frac{g!}{\alpha^{ND - |\Delta_j|}} \int_{|x_\mu - x_\nu|>\alpha} \bra*{\mathcal{O}_{j_1}(x_1) \dots \mathcal{O}_{j_N}(x_N)}\,\d^Dx_1 \dots \d^Dx_N.
\end{align}
Using $\frac{\d}{\d l} \alpha(l) = \alpha(l)$ we then find
\begin{align}
\frac{\d}{\d l}Z_N &= \sum_{j} \klb*{ \frac{\frac{\d}{\d l} g!}{\alpha^{ND- |\Delta_{j}|}} + \kla*{|\Delta_j| - ND} \frac{g!}{\alpha^{ND - |\Delta_j|}} }I_N + \sum_j \frac{g!}{\alpha^{ND-|\Delta_j|}} \frac{\d }{\d l} I_N(j).
\end{align}
Note that the only $l$-dependence of the integral is contained in the boundary conditions $|x_\mu-x_\nu|>\alpha(l)$ for $\mu,\nu = 1,\dots,N$ with $\mu<\nu$.
For the derivative it is therefore sufficient to only consider
\begin{align}
\frac{\d}{\d l}\int_{|x_\mu - x_\nu|>\alpha} &= \frac{\d}{\d l} \int \prod_{\mu<\nu} \Theta\kla*{|x_\mu-x_\nu|-\alpha} \\
&= -\alpha \int \sum_{\mu'<\nu'} \delta\kla*{|x_{\mu'}-x_{\nu'}| - \alpha} \prod_{ \mu<\nu , (\mu,\nu)\neq (\mu',\nu')} \Theta\kla*{|x_\mu-x_\nu|-\alpha}\\
&= -\alpha \sum_{\mu'<\nu'} \int_{|x_\mu - x_\nu|>\alpha,(\mu,\nu)\neq (\mu',\nu')} \delta\kla*{|x_{\mu'}-x_{\nu'}| - \alpha}.
\end{align}
After substituting $x_{\mu'}=x+x_{\nu'}$ and switching to $D$-dimensional spherical coordinates for the integration on $x$ this expression simplifies to
\begin{align}
-\alpha \sum_{\mu'<\nu'} \int_{|x_\mu - x_\nu|>\alpha,(\mu,\nu)\neq (\mu',\nu')} \delta\kla*{|x| - \alpha} S_D|x|^{D-1}
\end{align}
where we made the assumption that the angular dependence of the integrand is negligible for the small magnitude $|x|=\alpha$.
The boundary conditions are either $|x_\mu-x_\nu|>\alpha$ for $\mu,\nu \neq \mu'$ or $|x+x_{\nu'}-x_\nu|>\alpha$ for $\mu=\mu'$ and $\nu\neq \nu'$.
Here we will make the approximation to neglect $x$ in the second type of boundary condition.
On the one hand this is a bit brutal because $|x|=\alpha$ and thus ignoring this term when the RHS is of identical magnitude might seem dangerous.
On the other hand we know that averaged over the angles $x=0$ and also $\alpha$ only serves as a minimal distance between the coordinates.
Essentially for most of the parameter space neglecting a shift of $x$ does not matter.
However, the more honest motivation for this approximation is that it leads to a very simple structure that we can work with; potentially a better treatment here could even lead so slightly modified flow equations, but we will not consider this at any higher level of rigour.\\
Let us for now tentatively accept this and note that it leads to
\begin{align}
\frac{\d}{\d l}\int_{|x_\mu - x_\nu|>\alpha} &\approx -S_D\alpha^D \sum_{\mu'<\nu'} \int_{|x_\mu - x_\nu|>\alpha,\mu\neq \mu'}.
\end{align}
Inserting back into our integral gives
\begin{align}
\frac{\d}{\d l} I_N(j) &= -S_D\alpha^D \sum_{\mu'<\nu'} \int \bra*{\mathcal{O}_{j_1}(x_1) \dots \mathcal{O}_{j_{\mu'}}(x_{\nu'}+\alpha)\dots \mathcal{O}_{j_{\nu'}}(x_{\nu'}) \dots \mathcal{O}_{j_{N}}(x_{N})}\\
&=-S_D\alpha^D \sum_{\mu'<\nu'} \sum_{j'} \frac{C_{j_{\mu'}, j_{\nu'}, j'}}{\alpha^{\Delta_{j_{\mu'}} + \Delta_{j_{\nu'}} - \Delta_{j'}}} \int \bra*{\mathcal{O}_{j_1}(x_1) \dots \mathcal{O}_{j'}(x_{\nu'}) \dots \mathcal{O}_{j_{N}}(x_{N})}\\
&=-S_D\alpha^D \sum_{\mu'<\nu'} \sum_{j'} \frac{C_{j_{\mu'}, j_{\nu'}, j'}}{\alpha^{\Delta_{j_{\mu'}} + \Delta_{j_{\nu'}} - \Delta_{j'}}} I_{N-1}(j_1,\dots, j',\dots, j_N). \label{eqn:1loop shell integral OPE}
\end{align}
Here we made use of the operator product expansion
\begin{align}
\mathcal{O}_j(x_1) \mathcal{O}_k(x_2) &\overset{x_1\rightarrow x_2}{\sim} \sum_m \frac{C_{j,k,m}}{|x_1-x_2|^{\Delta_{j} + \Delta_{k} - \Delta_{m}}} \mathcal{O}_m\kla*{\frac{x_1+x_2}{2}}
\end{align}
where the $C$'s are the so-called \emph{structure coefficients} of the OPE.
One should be somewhat careful with this identity in the sense that there are a few implicit assumptions behind this asymptotic limit.
It is generally only valid inside an expectation value and it may generate operators not previously contained in the set we started with; this is the sense in which one usually assumes the $\mathcal{O}_j$ to be complete, i.e. all of their OPE's do not lead to any new terms.
The former is more of a mathematical subtlety that is presumably safe to ignore outside special cases while the latter is potentially a problem for the RG as this may generate entire families of operators that quickly make any high order calculation completely infeasible.
An argument that is somewhat symptomatic of RG is therefore to use the tree-level scaling as a justification for throwing away highly irrelevant terms that appear in this way without looking at their contributions at higher orders.\\
Plugging our result back into the derivative of the $N$-term partition function yields
\begin{align}
\frac{\d}{\d l} Z_N &= \bsplit{&\sum_j \frac{g!}{\alpha^{ND-|\Delta_j|}} \klb*{ \abs*{\frac{g'}{g}} + |\Delta_j| - ND }I_N(j) \\
&- \sum_j \frac{g!}{\alpha^{ND - |\Delta_j|}} S_D\alpha^D \sum_{\mu'<\nu'} \sum_{j'} \frac{C_{j_{\mu'}, j_{\nu'}, j'}}{\alpha^{\Delta_{j_{\mu'}} + \Delta_{j_{\nu'}} - \Delta_{j'}}} I_{N-1}(j',j\setminus \klc*{j_{\mu'}, j_{\nu'}}).}
\end{align}
Let us write out the condition $\frac{\d}{\d l} \frac{Z}{Z_0} = 0$ explicitly to better understand the structure.
To order $N=3$ one finds
\begin{align}
0 &= \bsplit{&\sum_{j_1} \frac{g_{j_1}}{\alpha^{D-\Delta_{j_1}}} \klb*{\frac{g_{j_1}'}{g_{j_1}} + \Delta_{j_1} - D} I_1(j_1)\\
& + \frac{1}{2}\sum_{j_1,j_2} \frac{g_{j_1}g_{j_2}}{\alpha^{2D-\Delta_{j_1} - \Delta_{j_2}}} \klb*{\frac{g_{j_1}'}{g_{j_1}} + \frac{g_{j_2}'}{g_{j_2}} + \Delta_{j_1} + \Delta_{j_2} - 2D} I_2(j_1,j_2)\\
&\qquad - \frac{1}{2}\sum_{j_1,j_2} \frac{g_{j_1}g_{j_2}}{\alpha^{2D-\Delta_{j_1} - \Delta_{j_2}}} S_D\alpha^D \sum_{j'} \frac{C_{j_{1}, j_{2}, j'}}{\alpha^{\Delta_{j_{1}} + \Delta_{j_{2}} - \Delta_{j'}}} I_{1}(j')\\
& + \frac{1}{6}\sum_{j_1,j_2,j_3} \frac{g_{j_1}g_{j_2}g_{j_3}}{\alpha^{3D-\Delta_{j_1} - \Delta_{j_2} - \Delta_{j_3}}} \klb*{\frac{g_{j_1}'}{g_{j_1}} + \frac{g_{j_2}'}{g_{j_2}} + \frac{g_{j_3}'}{g_{j_3}} + \Delta_{j_1}+ \Delta_{j_2} + \Delta_{j_3} - 3D} I_3(j_1,j_2,j_3)\\
&\qquad - \frac{1}{6}\sum_{j_1,j_2,j_3} \frac{g_{j_1}g_{j_2}g_{j_3}}{\alpha^{3D-\Delta_{j_1} - \Delta_{j_2} - \Delta_{j_3}}} S_D\alpha^D \sum_{j'} \sum_{\mu'=1}^2 \sum_{\nu'=\mu'+1}^3 \frac{C_{j_{\mu'}, j_{\nu'}, j'}}{\alpha^{\Delta_{j_{\mu'}} + \Delta_{j_{\nu'}} - \Delta_{j'}}} I_{2}(j',j\setminus \klc*{j_{\mu'}, j_{\nu'}}).}
\end{align}
If we go through these terms order by order we first find the condition
\begin{align}
\frac{g_{j}}{\alpha^{D-\Delta_{j}}} \klb*{\frac{g_{j}'}{g_{j}} + \Delta_{j} - D} &=  \frac{1}{2} \sum_{j_1,j_2} \frac{g_{j_1}g_{j_2}}{\alpha^{2D-\Delta_{j_1} - \Delta_{j_2}}} S_D\alpha^D \frac{C_{j_{1}, j_{2}, j}}{\alpha^{\Delta_{j_{1}} + \Delta_{j_{2}} - \Delta_{j}}}
\end{align}
which can be rewritten as
\begin{align}
g_j' &= (D-\Delta_j)g_j + \frac{S_D}{2} \sum_{j_1,j_2} C_{j_1,j_2,j} g_{j_1} g_{j_2}.
\end{align}
Note the summations on $j_1$ and $j_2$ only run over the perturbations while the coupling $g_j$ corresponding to an operator $\mathcal{O}_j$ may also be present in the fixed point action $S_0$.\\
The second order terms demand that
\begin{align}
\sum_{j_1,j_2} &\frac{g_{j_1}g_{j_2}}{\alpha^{2D-\Delta_{j_1} - \Delta_{j_2}}} \klb*{\frac{g_{j_1}'}{g_{j_1}} + \frac{g_{j_2}'}{g_{j_2}} + \Delta_{j_1} + \Delta_{j_2} - 2D} I_2(j_1,j_2) \\
&= \frac{1}{3}\sum_{j_1,j_2,j_3} \frac{g_{j_1}g_{j_2}g_{j_3}}{\alpha^{3D-\Delta_{j_1} - \Delta_{j_2} - \Delta_{j_3}}} S_D\alpha^D \sum_{j'} \bsplit{&\frac{C_{j_{1}, j_{2}, j'}}{\alpha^{\Delta_{j_{1}} + \Delta_{j_{2}} - \Delta_{j'}}} I_{2}(j',j_3)\\
&+\frac{C_{j_{1}, j_{3}, j'}}{\alpha^{\Delta_{j_{1}} + \Delta_{j_{3}} - \Delta_{j'}}} I_{2}(j',j_2)\\
&+ \frac{C_{j_{2}, j_{3}, j'}}{\alpha^{\Delta_{j_{2}} + \Delta_{j_{3}} - \Delta_{j'}}} I_{2}(j',j_1)}\\
&= \sum_{j_1',j_2',j_3'} \frac{g_{j_1'}g_{j_2'}g_{j_3'}}{\alpha^{3D-\Delta_{j_1'} - \Delta_{j_2'} - \Delta_{j_3'}}} S_D\alpha^D \sum_{j'}\frac{C_{j_{1}', j_{2}', j'}}{\alpha^{\Delta_{j_{1}'} + \Delta_{j_{2}'} - \Delta_{j'}}} I_{2}(j',j_3')
\end{align}
i.e. $(j_1,j_2)=(j',j_3')$ and thus upon symmetrizing
\begin{align}
\frac{g_{j_1}g_{j_2}}{\alpha^{2D-\Delta_{j_1} - \Delta_{j_2}}} \klb*{\frac{g_{j_1}'}{g_{j_1}} + \frac{g_{j_2}'}{g_{j_2}} + \Delta_{j_1} + \Delta_{j_2} - 2D} &= \frac{1}{2}\sum_{j_1',j_2'}\bsplit{& \frac{g_{j_1'}g_{j_2'}g_{j_2}}{\alpha^{3D-\Delta_{j_1'} - \Delta_{j_2'} - \Delta_{j_2}}} S_D\alpha^D \frac{C_{j_{1}', j_{2}', j_1}}{\alpha^{\Delta_{j_{1}'} + \Delta_{j_{2}'} - \Delta_{j_1}}}\\
&+ \frac{g_{j_1'}g_{j_2'}g_{j_1}}{\alpha^{3D-\Delta_{j_1'} - \Delta_{j_2'} - \Delta_{j_1}}} S_D\alpha^D \frac{C_{j_{1}', j_{2}', j_2}}{\alpha^{\Delta_{j_{1}'} + \Delta_{j_{2}'} - \Delta_{j_2}}}}
\end{align}
which then reduces to
\begin{align}
g_{j_1}'g_{j_2} + g_{j_2}'g_{j_1} &= (2D-\Delta_{j_1} - \Delta_{j_2})g_{j_1} g_{j_2} + \frac{S_D}{2} \sum_{j_1',j_2'} \klb*{C_{j_1',j_2',j_1} g_{j_1'} g_{j_2'} g_2 + C_{j_1',j_2',j_2} g_{j_1'} g_{j_2'} g_1}.
\end{align}
This is not equivalent to the much simpler condition found at first order, but it nevertheless automatically satisfied if the first condition holds.
Thus we do not find any modification at this level and although it is a combinatoric nightmare the author is convinced that this pattern continues to arbitrary order.
As such the RG approach presented here appears to be insufficient to generate corrections beyond the 1-loop order -- which may be perfectly sufficient for our use case -- but still makes for an interesting observation.
\TODO{It is certainly possible that something was missed that would break this argument. Perhaps the brutal approximation with the boundary conditions is at fault, since that only starts to become relevant for $N\ge 3$.
Additionally there might be some subtlety to the condition $\frac{\d}{\d l} \frac{Z}{Z_0} = 0$ that only allows for first order corrections; maybe a more careful approach would be required.}\\
To this one-loop level we thus obtain a structurally identical action where the new couplings are
\begin{align}
\frac{\d g_j}{\d l} &= (D - \Delta_j)g_j + \frac{1}{2} S_D \sum_{k,m} C_{k,m,j} g_k g_m.\label{eqn:beta function 1loop}
\end{align}
Thus all one has to do to this level is to identify the scaling dimensions of all operators in the theory and then compute the structure coefficients of the OPE.\\
Let us finish with a rather subtle but important point.
While the set of operators only contains the perturbations, the OPE may also involve terms contained in $S_0$.
By construction the scaling dimension of anything contained in $S_0$ must always be given by $\Delta=D$ because otherwise $S_0$ would not correspond to a fixed point.
For a simple example say that $\mathcal{O}_0$ is the operator describing $S_0$.
Under the RG there can now be a change of $g_0$ due to a feedback of the OPE of the form $C_{j,k,0}$.
Note however, that even if $C_{0,0,0} \neq 0$ this would \emph{not} appear in beta function!
This may be somewhat confusing so let us elaborate on this point.
The original summation in the expansion of the partition function only contains the perturbations, thus $j,k,m,\dots\neq 0$.
The OPE however can also produce an $m=0$-term and thus the sum in \autoref{eqn:1loop shell integral OPE} generally runs over more indices.
The sums in the beta function in \autoref{eqn:beta function 1loop} once again come from the expansion though, so they \emph{are} restricted to $j,k\neq 0$.
\section{Model Hamiltonian}
The full model is given by the Hamiltonian $H=H_0+H_1+H_2 + H_3 + H_4$
\begin{align}
H_0 &= \frac{v}{2\pi} \intx{ }{ }{\klb*{K\kla*{\partial_x\theta(x)}^2 + \frac{1}{K} \kla*{\partial_x\phi(x)}^2}}{x},\\
H_1 &= \frac{vg_1}{\alpha} \sum_{s_a=\pm 1} \intx{ }{ }{\sin\kla*{2\phi(x)} \partial_x\phi(x+s_aa)}{x},\\
H_2 &= \sum_{s_a=\pm 1} \frac{vg_2}{\alpha^2} \intx{ }{ }{\cos\kla*{2\phi(x) - 2\phi(x+s_aa)}}{x},\\
H_3 &= \sum_{s_a=\pm 1} \frac{vg_3}{\alpha^2} \intx{ }{ }{\cos\kla*{2\phi(x) + 2\phi(x+s_aa)}}{x},\\
H_4 &= \sum_{s_a=\pm 1} vg_4 \intx{ }{ }{\partial_x\phi(x) \partial_x\phi(x+s_aa)}{x}.
\end{align}
The integrations are restricted to $[-L/2, L/2]$ and $a$ is a lattice spacing assumed to satisfy $a<\alpha$.\\
We work in $(1+1)$ dimensions and the field conjugate momentum $\Pi = \frac{1}{\pi} \partial_x \theta$ leads to $\partial_x\theta = \frac{\i}{Kv} \partial_\tau\phi$ where $\tau=\i t$ is the imaginary time.
The action formulation is thus $S=S_0 + S_1 + S_2 + S_3 + S_4$ where
\begin{align}
S_0 &= -\intx{0}{\beta}{\klb*{\Pi \i\partial_\tau\phi - H_0} }{\tau} = \frac{1}{2\pi K} \intx{0}{\beta}{\intx{ }{ }{\klb*{\frac{1}{v} \kla*{\partial_\tau \phi}^2 + v \kla*{\partial_x \phi}^2}}{x}}{\tau}
\end{align}
and $S_j = \intx{0}{\beta}{H_j}{\tau}$ for $j\in\klc*{1,2,3,4}$.
It is more convenient to express everything in terms of the euclidean coordinates $r = (x,y)$ with $y=v\tau$ which leads to
\begin{align}
S_0 &= \frac{1}{2\pi K} \int \kla*{\nabla\phi(r)}^2 \,\d^2r,\\
S_1 &= \frac{g_1}{\alpha} \sum_{s_a=\pm 1} \int \sin \kla*{2 \phi(r)} \partial_x \phi(r+s_aa) \,\d^2r,\\
S_2 &= \frac{g_2}{\alpha^2} \sum_{s_a=\pm 1} \int \cos\kla*{2\phi(r) - 2\phi(r+s_aa)}\,\d^2r,\\
S_3 &= \frac{g_3}{\alpha^2} \sum_{s_a=\pm 1} \int \cos\kla*{2\phi(r) + 2\phi(r+s_aa)}\,\d^2r,\\
S_4 &= g_4 \sum_{s_a=\pm 1} \int \partial_x\phi(r) \partial_x\phi(r+s_aa)\,\d^2r
\end{align}
where $g_1,g_2,g_3,g_4,K$ are dimensionless parameters and $\kla*{\nabla\phi}^2 = \kla*{\partial_x\phi}^2 + \kla*{\partial_y\phi}^2$.
Keep in mind that $\phi$ is a real bosonic scalar field given by a linear combination of bosonic operators $b_q$ and can therefore be split into two parts $\phi^{(\pm)}$ containing creation and annihilation operators respectively.
Their commutator satisfies
\begin{align}
\klb*{\phi^{(+)}(r), \phi^{(-)}(0)} &= \frac{K}{2} \log\klb*{\frac{2\pi}{L} \sqrt{\alpha^2 + r^2}}
\end{align}
and the fundamental correlation function is asymptotically given by
\begin{align}
\Phi(r) &= \bra*{\klb*{\phi(r) - \phi(0)}^2} = -2 \bra*{\kla*{\phi(r) - \phi(0)} \phi(0)} \overset{r\rightarrow 0}{\sim} \frac{K}{2} \log\kla*{\frac{r^2 + \alpha^2}{\alpha^2}}.
\end{align}
We also give the general formula for products of vertex operators
\begin{align}
\bra*{\prod \powe{\i A_j\phi(r_j)}} &= \powe{-\frac{1}{2} \sum A_j A_k \bra*{\phi(r_j-r_k) \phi(0)}} = \powe{-\frac{1}{2} \kla*{\sum A_j}^2 \bra*{\phi(0)^2}} \powe{\frac{1}{2} \sum_{j<k} A_jA_k \Phi(r_j-r_k)}.
\end{align}
Formally $\bra*{\phi(0)^2}=\infty$ for an infinite system giving rise to the neutrality condition $\sum_j A_j =0$.
We will see that minor violations of this condition via source terms can be dealt with by delaying the thermodynamic limit.
\subsection{Scaling dimensions and normalization}
Let us now analyse the operators present in our perturbation.
We will often make use of the formulas
\begin{align}
\lim_{J_1,J_2\rightarrow 0} \partial_{J_1} \partial_{J_2} \powe{A_{11}J_1^2 + A_{22}J_2^2 + A_{12} J_1J_2 + B_1J_1 + B_2J_2} &= A_{12} + B_1B_2,\\
\lim_{J_1,J_2,J_3,J_4\rightarrow 0} \partial_{J_1} \partial_{J_2} \partial_{J_3} \partial_{J_4} \powe{\sum_{j\le k} A_{jk} J_jJ_k} &= A_{12} A_{34} + A_{13}A_{24} + A_{14}A_{23}.
\end{align}
%J1,J2,J3,J4,A12,A13,A14,A23,A24,A34,B1,B2,B3,B4,A11,A22,A33,A44=sy.symbols("J_1 J_2 J_3 J_4 A_12 A_13 A_14 A_23 A_24 A_34 B_1 B_2 B_3 B_4 A_11 A_22 A_33 A_44")
%expr1=sy.exp(J1**2*A11 + J2**2*A22+J1*J2*A12+J1*B1+J2*B2)
%expr2=sy.exp(J1**2*A11 + J2**2*A22+ J3**2*A33+ J4**2*A44+J1*J2*A12+J1*J3*A13+J1*J4*A14+J2*J3*A23+J2*J4*A24+J3*J4*A34+J1*B1+J2*B2+J3*B3+J4*B4)
%expr1.diff(J1).diff(J2).subs(J1,0).subs(J2,0)
%expr2.diff(J1).diff(J2).diff(J3).diff(J4).subs(J1,0).subs(J2,0).subs(J3,0).subs(J4,0)
Also for convenience of notation limits will often by implicit.
The standard replacement rule for the shifted operators is
\begin{align}
\partial_x\phi(x+s_aa,y) &= \partial_{x_1} \phi(x_1,y) = -\i\partial_{x_1} \partial_{J_1} \powe{\i J_1 \phi(r_1)}
\end{align}
where $x_1=x+s_aa$ and $r_1=(x_1,y)$.\\
Let us start by considering
\begin{align}
\sum_{s_a,s_a'}&\bra*{\partial_x\phi(r) \partial_x \phi(r+s_aa) \partial_x\phi(r+\epsilon) \partial_x\phi(r+\epsilon + s_a'a)} \\
&= \sum_{x_2=x_1\pm a}\sum_{x_4=x_3\pm a} \partial_x^4 \partial_J^4 \bra*{ \powe{\i J_1\phi(r_1)} \powe{\i J_2\phi(r_2)} \powe{\i J_3\phi(r_3)} \powe{\i J_4\phi(r_4)}} \\
&= \sum \partial^4 \powe{-\frac{1}{2} \kla*{\sum J}^2 \bra*{\phi(0)^2}} \powe{\frac{1}{2} \sum_{j<k} J_jJ_k \Phi(r_j-r_k)}\\
&= \sum \partial_x^4 \frac{1}{2} \klb*{\Phi(r_1-r_2) \Phi(r_3-r_4) + \Phi(r_1-r_3) \Phi(r_2-r_4) + \Phi(r_2-r_3) \Phi(r_1-r_4)}\\
&= \frac{K}{2} \sum \partial_x^4 \klb*{\log|r_1-r_2|\log |r_3-r_4| + \log|r_1-r_3|\log |r_2-r_4| + \log|r_2-r_3|\log |r_1-r_4|}\\
&= \frac{K}{2} \sum \klb*{\frac{-1}{(r_1-r_2)^2} \frac{-1}{(r_3-r_4)^2} + \frac{-1}{(r_1-r_3)^2} \frac{-1}{(r_2-r_4)^2} + \frac{-1}{(r_2-r_3)^2} \frac{-1}{(r_1-r_4)^2}}\\
&= \frac{K}{2} \sum \klb*{ \frac{1}{a^2} \frac{1}{a^2} + \frac{1}{\epsilon^2} \frac{1}{\kla*{\epsilon+s_aa + s_a'a}^2} + \frac{1}{\kla*{\epsilon + s_aa}^2} \frac{1}{\kla*{\epsilon + s_a'a}^2}}.
\end{align}
A few details were left out here for efficiency.
First, we just write $\partial_x^4$ for $\partial_{x_1}\dots \partial_{x_4}$ and likewise the derivatives with respect to $J$ and the summations are condensed too.
Also, we automatically enumerate the arguments with $r_1,\dots,r_4$ from left to right so $r_1=(x,y)$, $r_2=(x+s_aa,y)$, $r_3=(x+\epsilon_x,y+\epsilon_y)$ and $r_4=(x+s_aa+\epsilon_x,y+\epsilon_y)$.
Additionally the terms $A_jk$ are of the form $\Phi(r_jk) - \bra*{\phi(0)^2}$ and we have dropped the constant since it vanishes upon differentiation anyway.
It is clear that this has to be the case for we would otherwise run into divergences when taking the thermodynamic limit.\\
The previous result still requires some work because we have to discuss the hierarchy of $\epsilon$ and $a$.
One could argue that we only ever use the expansion for $\epsilon \approx \alpha > a$ and thus should neglect $a$ in comparison to $\epsilon$.
%For the purpose of analysing the scaling dimension and structure coefficients we should really think of $\epsilon\rightarrow 0$ though.
%While this is a bit inconsistent with the later use it avoids a lot of ambiguity with interpreting the OPE itself.
%Essentially we could hardly read of the scaling dimension from a term like $\frac{1}{\epsilon^2} + \frac{1}{a^2}$ unless we assume that the dominant, singular part is the one containing $\epsilon$.
\QnA{This is a little bit concerning because this essentially means we consider $\epsilon\gg a$ while simultaneously ignoring this offset of $\epsilon$ in the derivation of the kinetic term of the Luttinger Liquid.
This justifies later neglecting the shift in the argument of $H_4$ but at the same time makes it rather awkward to claim that the shift by $a$ in all the sine-Gordon-like terms is somehow essential.
I do not see a way to avoid all of the problems though, so there seems to be some kind of minimal ``voodoo-magic'' part associated with the entire LL thing.
In defence of my approach it has to be said that it avoids quite a lot of other issues present -- but usually just ignored -- in other people's research. So I suppose I am at least not significantly more hand-wavy than everyone else.}\\
%With this approximation we can then reduce the expression to
%\begin{align}
%\frac{K}{2} \klb*{\frac{8}{a^4} + \frac{2}{\epsilon^2a^2} + \frac{2}{\epsilon^4}} \sim \frac{K}{\epsilon^4}.
%\end{align}
%Thus the scaling dimension of this operator is 2, which is not too surprising as it coincides with the simple approach of ``power-counting'', and the normalization is $\frac{1}{\sqrt{K}}$.
%The first type of operator we define is therefore
%\begin{align}
%P(r) &= \frac{1}{\sqrt{K}} \sum_{s_a=\pm 1} \partial_x\phi(r) \partial_x\phi(r+s_aa).
%\end{align}
With this approximation we can then reduce the expression to
\begin{align}
\frac{K}{2} \klb*{\frac{4}{a^4} + \frac{8}{\epsilon^4}}.
\end{align}
Thus the scaling dimension of this operator is 2, which is not too surprising as it coincides with the simple approach of ``power-counting'', and the normalization is $\frac{1}{\sqrt{4K}}$.
The first type of operator we define is therefore
\begin{align}
P(r) &= \frac{1}{\sqrt{4K}} \sum_{s_a=\pm 1} \partial_x\phi(r) \partial_x\phi(r+s_aa).
\end{align}
(\TODO{choice of hierarchy ends up being a factor of two here; check consequence for RG flow later})\\
Next we consider
\begin{align}
\sum_{s_a,s_a'} &\bra*{\powe{\i n\phi(r)} \partial_x\phi(r+s_aa) \powe{\i m\phi(r+\epsilon)} \partial_x\phi(r+s_a'a+\epsilon)}\\
&= \sum \partial^2 \bra*{\powe{\i n\phi(r_1)}\powe{\i J_1\phi(r_2)}\powe{\i m\phi(r_3)}\powe{\i J_2\phi(r_4)}}\\
&= \sum \partial^2 \powe{\frac{1}{2} \klb*{nJ_1 \Phi_{12} + nm \Phi_{13} + nJ_2 \Phi_{14} + J_1m\Phi_{23} + J_1J_2 \Phi_{24} + mJ_2 \Phi_{34}}} \delta_{n,-m}\\
&= \sum \powe{\frac{1}{2} nm \Phi_{13}} \partial_x^2 \klb*{\Phi_{24} + \kla*{n\Phi_{12} + m\Phi_{23}} \kla*{m\Phi_{34} + n\Phi_{14}}}\\
&= K \powe{-\frac{1}{2}n^2 \Phi(\epsilon)} \sum \partial_x^2 \klb*{\log |r_2-r_4| + Kn^2 \kla*{\log |r_1-r_2| - \log |r_2-r_3|} \kla*{\log|r_1-r_4| - \log|r_3-r_4|}}\\
&= K \powe{-\frac{1}{2}n^2 \Phi(\epsilon)} \sum \klb*{ \frac{-1}{(r_2-r_4)^2} + Kn^2 \kla*{\frac{\sgn(x_2-x_1)}{|r_2-r_1|} -\frac{\sgn(x_2-x_3)}{|r_2-r_3|}} \kla*{\frac{\sgn(x_4-x_3)}{|r_4-r_3|} - \frac{\sgn(x_4-x_1)}{|r_4-r_1|}}}.
\end{align}
Once again this is as far as we can go without making a choice for the hierarchy.
As before we will take the practical point of view and put $\epsilon>\alpha>a$ leading to
\begin{align}
K \kla*{\frac{\alpha}{\epsilon}}^{\frac{K}{2}n^2} \klb*{\frac{-4}{\epsilon^2} + Kn^2 \frac{2}{\epsilon} \frac{-2}{\epsilon}} = -4K (1+Kn^2) \frac{1}{\epsilon^{\frac{K}{2}n^2+2}} \alpha^{\frac{K}{2}n^2}.
\end{align}
(\TODO{choice of hierarchy is 1 vs. $1+Kn^2$})\\
The scaling dimension is $\Delta^W_n = \frac{K}{4}n^2+1$ and we thus define the operators
\begin{align}
W_n(r) &= \frac{1}{2\i} \frac{1}{\sqrt{K \kla*{1 + Kn^2}}} \frac{1}{\alpha^{\Delta^W_n-1}} \sum_{s_a=\pm 1} \powe{\i n\phi(r)} \partial_x\phi(r+s_aa).
\end{align}
Note that the factor of $\frac{1}{2\i}$ corresponds to the complex representation of $\sin(2\phi)$.\\
Note that using the commutation relations laid out before we have
\begin{align}
\powe{\i n\phi(r) + \i ns\phi(r')} &= \powe{\i n\phi(r)}\powe{\i ns\phi(r')} \powe{-\frac{1}{2} \klb*{\i n\phi(r), \i ns\phi(r')}} \\
&= \powe{\i n\phi(r)}\powe{\i ns\phi(r')} \powe{\frac{n^2s}{2} \klb*{\phi^{(+)}(r) + \phi^{(-)}(r), \phi^{(+)}(r') + \phi^{(-)}(r')}}\\
&= \powe{\i n\phi(r)}\powe{\i ns\phi(r')} \powe{\frac{n^2s}{2} \klb*{\phi^{(-)}(r), \phi^{(+)}(r')} + \klb*{\phi^{(+)}(r), \phi^{(-)}(r')}} = \powe{\i n\phi(r)}\powe{\i ns\phi(r')}.
\end{align}
Next we consider 
\begin{align}
\sum_{s_a,s_a'}& \bra*{\powe{\i n\phi(r) + \i n s\phi(r+s_aa)}\powe{\i m\phi(r+\epsilon) + \i m s \phi(r+\epsilon+s_a'a)}}\\
&= \delta_{n(1+s),-m(1+s)} \sum \powe{\frac{1}{2} \klb*{n^2s\Phi_{12} + nm\Phi_{13} + nms \Phi_{14} + nsm\Phi_{23} + nsms\Phi_{24} + m^2s\Phi_{34}}}.
\end{align}
Once again let $\epsilon>\alpha>a$ and rewrite as
\begin{align}
4\delta_{n(1+s),-m(1+s)} \kla*{\frac{\epsilon}{\alpha}}^{\frac{K}{2}\klb*{nm(1+s)^2}}.
\end{align}
Evidently this can only be correct for $s=+1$.
Otherwise we should go back to the general result and expand
\begin{align}
\sum &\powe{-\frac{1}{2} \klb*{n^2\Phi_{12} - nm\Phi_{13} + nm \Phi_{14} + nm\Phi_{23} - nm\Phi_{24} + m^2\Phi_{34}}}\\
&\sim \sum \powe{\frac{K}{2}nm \klb*{\log|\frac{r_1-r_3}{r_1-r_4}| + \log|\frac{r_2-r_3}{r_2-r_4}|}} = \sum \powe{\frac{K}{2}nm \klb*{\log|\frac{\epsilon}{\epsilon + s_a'a}| + \log|\frac{\epsilon + s_aa}{\epsilon}|}}\\
&\sim \sum \abs*{\frac{1 + s_aa/\epsilon}{1 + s_a'a/\epsilon}}^{\frac{K}{2}nm} \sim \sum \abs*{ 1 + (s_a - s_a') \frac{a}{\epsilon} + \frac{a^2}{\epsilon^2} }^{\frac{K}{2}nm} (???)
\end{align}
\QnA{Maybe I can just set $a=0$ for the RG? That would just make everything much less horrible looking and probably also avoids getting lost in unnecessary details.
Then again we have very good arguments \emph{against} $a=0$ as this would for example eliminate the $H_1$-term.
Of course we could just take the stance that we know that this is not quite a total derivative and thus argue the problem away; the RG frankly does not seem to care one bit about this particular point.}
\subsection{OPE}



\end{document}